% from ega1-6-fr.tex
\oldpage[I]{151}
    égaux à un même élément $h\in B-\mathfrak{j}_x$~; on en conclut que
    $P_j(hb_1,\ldots,hb_n,hg)=0$ pour $1\leq j\leq m$. Cela étant, considérons
    l'homomorphisme $\rho$ de $A[X_1,\ldots,X_n]$ dans l'anneau de fractions
    $B_{hg}$ qui applique $X_i$ sur $hb_i/hg$ ($1\leq i\leq n$)~; l'image de
    $\mathfrak{a}$ par cet homomorphisme est 0, et il en est \emph{a fortiori}
    de même de l'image par $\rho$ du noyau $\alpha^{-1}(0)$. Donc $\rho$ se
    factorise en
    $A[X_1,\ldots,X_n]\xrightarrow{\alpha}C\xrightarrow{\gamma}B_{hg}$, avec
    $\gamma(c_i)=hb_i/hg$, et il est clair que si $i_x$ est l'homomorphisme
    canonique $B_{hg}\to B_x$, le diagramme
    \[
      \xymatrix{
        C\ar[r]^\gamma\ar[d]_{i_y} & B_{hg}\ar[d]^{i_x}\\
        C_y\ar[r]_\varphi & B_x
      }
      \tag{6.5.1.1}\label{I.6.5.1.1-fr}
    \]
    est commutatif~; on a donc $\varphi=\gamma_x$, et comme $\varphi$ est un
    homomorphisme local, ${}^a\gamma(x)=y$~; $f=({}^a\gamma,\widetilde{\gamma})$
    est donc un $S$-morphisme du voisinage $D(hg)$ de $x$ dans $Y$ qui répond à
    la question.
\end{enumerate}

\begin{corollary}[6.5.2]
\label{I.6.5.2-fr}
Sous les hypothèses de (\hyperref[I.6.5.1-fr]{6.5.1, (ii)}), si en outre $X$
est de type fini sur $S$, on peut supposer le morphisme $f$ de type fini.
\end{corollary}

Cela résulte de (\hyperref[I.6.3.6-fr]{6.3.6}).

\begin{corollary}[6.5.3]
\label{I.6.5.3-fr}
Supposons vérifiées les hypothèses de
(\hyperref[I.6.5.1-fr]{6.5.1, (ii)}), et supposons en outre que $Y$ soit
intègre, et $\varphi$ un homomorphisme injectif. Alors on peut supposer que
$f=({}^a\gamma,\widetilde{\gamma})$, où $\gamma$ est injectif.
\end{corollary}

En effet, on peut supposer $C$ intègre
(\hyperref[I.5.1.4-fr]{5.1.4}), donc $i_y$ injectif~; il résulte alors du
diagramme (\hyperref[I.6.5.1.1-fr]{6.5.1.1}) que $\gamma$ est injectif.

\begin{proposition}[6.5.4]
\label{I.6.5.4-fr}
Soient $f=(\psi,\theta):X\to Y$ un morphisme de type fini, $x$ un point de
$X$, $y=\psi(x)$.
\begin{enumerate}
  \item[(i)] Pour que $f$ soit une immersion locale au point $x$
    (\hyperref[I.4.5.1-fr]{4.5.1}), il faut et il suffit que
    $\theta_x^\sharp:\mathscr{O}_y\to\mathscr{O}_x$ soit surjectif.
  \item[(ii)] On suppose en outre $Y$ localement noethérien. Pour que $f$ soit
    un isomorphisme local au point $x$
    (\hyperref[I.4.5.2-fr]{4.5.2}), il faut et il suffit que
    $\theta_x^\sharp$ soit un isomorphisme.
\end{enumerate}
\end{proposition}

\begin{enumerate}
  \item[(ii)] En vertu de (\hyperref[I.6.5.1-fr]{6.5.1}), il existe alors un
    voisinage ouvert $V$ de $y$ et un morphisme $g:V\to X$ tels que $g\circ f$
    (resp. $f\circ g$) soit défini et coïncide avec l'identité dans un
    voisinage de $x$ (resp. $y$), d'où on tire aisément que $f$ est un
    isomorphisme local.
  \item[(i)] La question étant locale sur $X$ et $Y$, on peut supposer $X$ et
    $Y$ affines, d'anneaux respectifs $A$, $B$~; on a
    $f=({}^a\varphi,\widetilde{\varphi})$, où $\varphi$ est un homomorphisme
    d'anneaux $B\to A$ qui fait de $A$ une $B$-algèbre de type fini~; on a
    $\varphi^{-1}(\mathfrak{j}_x)=\mathfrak{j}_y$ et l'homomorphisme
    $\varphi_x:B_y\to A_x$ déduit de $\varphi$ est surjectif. Soit $(t_i)$
    ($1\leq i\leq n$) un système de générateurs de la $B$-algèbre $A$~;
    l'hypothèse sur $\varphi_x$ implique qu'il existe des $b_i\in B$ et un
    $c\in B-\mathfrak{j}_y$ tels que, dans l'anneau de fractions $A_x$, on ait
    $t_i/1=\varphi(b_i)/\varphi(c)$ pour $1\leq i\leq n$. Par suite
    (\hyperref[I.1.3.3-fr]{1.3.3}), il existe $a\in A-\mathfrak{j}_x$ tel que,
    si l'on pose $g=a\varphi(c)$, on ait aussi
    $t_i/1=a\varphi(b_i)/g$ \emph{dans l'anneau de fractions $A_g$}. Cela étant,
    il existe par hypothèse un polynôme $Q(X_1,\ldots,X_n)$ à coefficients dans
    l'anneau $\varphi(B)$,
